Erklären sie den Stakeholder-Ansatz vs. Shareholder Ansatz

Im Stakeholder-Ansatz erfolgt eine Diskussion zwischen den
Beteiligten und eine gerechte Verteilung. Dieser scheitert an
Interessengegensätzen.

Die Stakeholder setzen die Unternehmensleitungen
zunehmend unter Druck. Diese sollen nicht nur ökonomisch,
sondern auch ökologisch und sozial handeln. 

Der Shareholder-Ansatz hingegen entspricht
marktwirtschaftlichem Wirtschaften

Die Shareholder wollen ihre Gewinne maximieren und
schließen mit allen Stakeholdern Verträge zu
Marktkonditionen (Arbeitsvertrag, Liefervertrag,
Versicherungsvertag etc.).

Shareholder beurteilen die ökologischen und sozialen
Aspekte unter einer Kosten/Nutzen Analyse eher nachrangig. 


Nennen sie einige der Stakeholder und Bestandteile des externen Rechnungswesen, welche spezifisch für Versicherungsunternehmen sind

Externes Rechnungswesen von Versicherungsunternehmen
- Vermögenslage 
- Finanzlage 
- Ertragslage

Addressaten und Ihre Interessen sind
- Rückversicherer
  - Kann die Rückversicherungsbeziehung langfristig fortgeführt werden?
- Versicherungsnehmer:innen (eigentlich nur besonderer Name für Kunden)
  - Wird die Versicherungsleistung im Fall der
    Fälle anstandslos gezahlt oder wird in einem
    Rechtsstreit um jeden Euro gestritten?
  - Wie hoch sind die Überschüsse, die meiner
    Rentenversicherung gutgeschrieben werden?
- Vermittler:innen 
  - Sind die Produkte für Vermittler attraktiv (Provision)
- Ratinggesellschaften
- Versicherungsaufsicht
- Eigenkapitalgeber:innen
  - Wie hoch ist die Dividende/Gewinnbeteiligung/Boni/Tantieme?
- Fiskus

Nennen sie Informationsquellen der Stakeholder für Versicherungsunternehmen 

- IFRS-Abschluss (International Financial Reporting Standards)
- Handelsbilanz
- Steuerbilanz
- Anhang
- Lagebericht
- Nachhaltigkeitsbericht

- ORSA (Own risk and solvency assessment)
- Solvabilitätsübersicht
- SFCR (Solvency and Financial Condition Reports, Solvency II)
- RSR (Regular Supervisory Report, Solvency II) 
- QRT (Quantitative Reporting Templates, Solvency II)


Nennen Sie die Aufgaben der wichtigsten Aufsichtsbehörden für Versicherungsunternehmen

- EU-Kommission, erlässt und überarbeitet Richtlinien (z.B. Solvency II Richtlinie (EURichtlinie 2009/138/EG), (Omnibus II))))
- EIOPA (European Insurance and Occupational Pensions Authority) erlässt für Solvency II bindende technische
  Standards
- Nationaler Gesetzgeber (local National Competent Authority (NCA)) implementiert diese Richtlinien, e.g. in Deutschland das VAG

Und die BaFin (Bundesanstalt für Finanzdienstleistungsaufsicht)
Die BaFin überwacht das deutsche Finanzsystem und bekommt daher die gesonderten Aufgaben zur:
• Erteilung Erlaubnis zum Geschäftsbetrieb
• Überwachung des Geschäftsbetriebs des VU (Einhaltung
  Gesetze, Wahrung der Belange der Versicherungsnehmer)
• Für Solvency II erlässt die BaFin
  Auslegungsentscheidungen
• Überwachung der dauernden Erfüllbarkeit der Verpflichtungen
  aus Versicherungsverträgen, der Solvabilität und der
  Risikotragfähigkeit (Solvenzaufsicht (Zahlungsfähigkeit von VU/KI))
• Überwachung einer ordnungsgemäßen Geschäftsorganisation
• Aufsichtliches Überprüfungsverfahren hinsichtlich Überprüfung
  der Strategie, Prozesse und Meldeverfahren, die sicherstellen,
  dass das VU alle Rechts-und Verwaltungsvorschriften einhält


Nennen sie die Besonderheiten in Bilanz und GuV von Versicherungsunternehmen

Aufgrund der besonderen Strukturierung des Haushalts von VU 
existiert neben dem HBG noch die Versicherungsunternehmens-Rechnungslegungsverordnung 
(RechVersV).

VU verfügen über kein großes
Sachanlagevermögen, daher wird in der
Bilanz – im Gegensatz zu Industrieunternehmen – nicht in Anlage- und
Umlaufvermögen unterschieden

Bei den Aktiva handelt es sich im
Wesentlichen um Kapitalanlagen, die zur
Bedeckung der versicherungstechnischen
Verpflichtungen dienen. Daher bietet sich eine
tiefere Gliederung nach einzelnen
Anlagearten an

Einen Großteil der Passiva stellen die
versicherungstechnischen Rückstellungen
dar. Sie werden daher separat ausgewiesen.

Gesamt-und Umsatzkostenverfahren nach
§275 HGB sind nicht geeignet, da
versicherungsspezifische Posten nicht
aufgeführt werden. Daher Umsatzkostenverfahren nach RechVersV verwendet.

Untergliederung in versicherungstechnische
und nichtversicherungstechnische
Rechnung erforderlich (Teilerfolgsprinzip)

• Angaben zu einzelnen Sparten im Anhang (Spartenrechnungsprinzip)
• Darstellung der Kosten innerhalb der GuV
  nach Funktionsbereichen (§43 RechVersV),
  keine Primärkosten, sondern
  Sekundärkostenprinzip
  (Verteilung der Kosten)

Zeitraumbezogenheit, Stochastizität, Kollektivbezogenheit, Immaterialität.

Versicherungsgeschäft ist durch die Vorauszahlung von Beiträgen und nachgelagerter Zahlung von Schäden/Rückkaufswerten
geprägt. Die Versicherer haben somit ein Anlageproblem. Industrieunternehmen haben hingegen ein Finanzierungsproblem, da vor
Leistungserstellung Produktionsanlagen errichtet werden müssen. Dies ermöglicht eine Vielfalt von Synergien, ein bekanntes Beispiel
ist der Erfolg von Berkshire Hathaway.

Geben sie Beispiel von Funktionsbereichen für die GuV nach (§43 RechVersV)

• Abschluss von Versicherungsverträgen
• Verwaltung von Versicherungsverträgen
• Verwaltung von Kapitalanlagen
• Aufwendungen für die Regulierung von Versicherungsfällen
• Aufwendungen, die nicht einem der o.g. Funktionsbereiche zugeordnet werden können


Geben sie Beispiel von Posten aus der Bilanz von Versicherungsunternehmen

Aktivseite: Keine Untergliederung in Anlage-und Umlaufvermögen im Ausweis
• Passivseite: Fremdkapital spielt keine Rolle (Verbot der FK-Aufnahme)


SLIDE 19

Geben sie Beispiel von Posten aus der GuV von Versicherungsunternehmen (vgl. Personen/Komposit VU)

SLIDE 18

Primär-und Sekundärkostenprinzip (Sammlung je Kostenart + Verteilung auf Funktionsbereiche)
• Unterteilung in Versicherungstechnik und Nicht-Versicherungstechnik
• Bei Sach-VU sind die Kapitalanlagen der Nichtversicherungstechnik zugeordnet, weil die Versicherungsnehmer nicht daran beteiligt sind.


Was ist das modifizierte Nettoprinzip von Versicherungen

Im Folgenden
BE = Beitragseinnahmen
RV = Rückversicherungsanteil

Bruttoprinzip beschreibt den unsaldierten Ausweis von Aktiva und Passiva in der Bilanz bzw. von
Erträgen und Aufwendungen in der GuV hinsichtlich des in Rückdeckung gegebenen
Versicherungsgeschäfts. Beim modifizierten Nettoprinzip wird vom Bruttobetrag in der Vorspalte
der Anteil der Rückversicherer abgezogen und dann in der Hauptspalte der Nettowert ausgewiesen.


INSERT IMAGE 20


Nettoprinzip beschreibt den saldierten Ausweis von Aktiva und Passiva in der Bilanz bzw. von
Erträgen und Aufwendungen in der GuV hinsichtlich des in Rückdeckung gegebenen
Versicherungsgeschäfts. Bruttogrößen werden mit den Anteilen des Rückversicherers saldiert.

INSERT IMAGE 21


Nenne sie wesentliche Kapitalanlagekategorien von Versicherungen und deren durchschnittlicher Anteil wenn möglich

Festverzinsliche Wertpapiere (80% >)
- Inhaberschuldverschreibungen (Anleihen)
- Namensschuldverschreibungen (nicht öffentlich handelbare Kredite, Ausgabe häufig durch Hypothekenbanken)
- Vorauszahlungen auf Versicherungsscheine (V können Darlehen aufnehmen und als Sicherheit angesparte Versicherungsansprüche hinterlegen)

Immobilien (5%)
- Hypotheken

Equities (5%)
- Aktien 
- Anteile an Investmentvermögen (Fonds)

Sowie die abgetrennt zu betrachten
- Beteiligungen (Mehr als 20% Besitz)
- Verbundene Unternehmen (Mehr als 50% Besitz)
